\section{The question that starts the project}\label{ch18:sec:question}
The contraction theorem of Chapter~\ref{ch:contraction} is about exact updates: each new point is $x_{n+1}=T(x_n)$, computed without any error. Real computations are rarely that clean. A computer keeps only finitely many digits, so it has to round each value of $T$ before storing it; a calculator that keeps two decimal places might store $0.87$ where the exact value is $0.875$. A measuring instrument has limited precision, and a value produced by a separate numerical procedure is itself only approximate. In each case the new point is close to $T$ of the old point, but not equal to it. So we ask: does the iteration still approach the fixed point? And if it does not, can we at least say how far from the fixed point it can end up?

This is a good question for a first investigation. It changes one clearly identified feature of a theorem we already understand and keeps everything else, so we do not have to invent a subject from nothing. This chapter carries out the whole investigation as a worked example, from a first guess to a proved result whose limits are clearly stated. The bounds we reach are standard consequences of the contraction theorem, not new mathematics. What the chapter shows is how an investigation proceeds.

\subsection*{The setting}
Throughout the chapter, $(X,d)$ is a nonempty complete metric space, and $T:X\to X$ is a contraction with factor $q$, where $0<q<1$:
\[
 d(Tx,Ty)\le q\,d(x,y)\qquad\text{for all }x,y\in X.
\]
Asking for $q>0$ costs nothing. A contraction with factor $0$ satisfies $d(Tx,Ty)\le0\le\frac12d(x,y)$, so it is also a contraction with factor $1/2$. By Theorem~\ref{thm:banach}, $T$ has exactly one fixed point, which we call $x_*$.

In place of the exact iterates we now have points $y_0,y_1,y_2,\ldots$ of $X$ and nonnegative numbers $\eps_0,\eps_1,\eps_2,\ldots$ such that
\[
 d(y_{n+1},T(y_n))\le\eps_n\qquad\text{for every }n\in\NN.
\]
The starting point $y_0$ may be any point of $X$. The point $y_{n+1}$ is what the update actually produced when it should have produced $T(y_n)$, and the number $\eps_n$ bounds the error made in that update. For short, we call $\eps_n$ the \emph{update error} at step~$n$; strictly speaking it is an upper bound for that error, and an upper bound is all we will need. If every $\eps_n$ is zero, then $y_{n+1}=T(y_n)$ for every $n$, and we are back to the exact iteration of Chapter~\ref{ch:contraction}.

Apart from these bounds, we assume nothing about the errors. The bounds $\eps_n$ may change from step to step, the errors need not follow any pattern, and they are not assumed to be random.

One more assumption is easy to overlook: every $y_n$ must be a point of $X$. An error bound alone cannot guarantee this. Suppose, for instance, that $X=[0,1]$ and a computed value comes out as $-0.001$. This number is very close to the interval, but it is not a point of $X$. The distance $d(y_{n+1},T(y_n))$ is then not even defined, because $d$ only measures distances between points of $X$, and $T$ cannot be applied at the next step. In an application, the requirement $y_n\in X$ has to be checked separately, just like the self-map condition of Section~\ref{ch12:sec:iteration}.
